\documentclass[../../main.tex]{subfiles}
\begin{document}

{\bf [解答]}

$\theta=\dfrac{2\pi}{5},\dfrac{4\pi}{5}$はいずれも$5\theta=2n\pi$をみたす($n\in\mathbb{Z}$)ので
\begin{align}
  &3\theta=2n\pi-2\theta \\
  \therefore
  &\cos3\theta=\cos2\theta
\end{align}
である．

以下$t=\cos\theta$とすれば倍角，三倍角の公式から
\begin{align}
  4t^3-3t=2t^2-1      \\
  (t-1)(4t^2+2t-1)=0
\end{align}
である．
明らかに$t \not 1$だから
\begin{align}
  4t^2+2t-1 = 0 \label{eq:1}
\end{align}
である．
ここで$\cos\left(\dfrac{2\pi}{5}\right)\not=\cos\left(\dfrac{4\pi}{5}\right)$より，$2$次方程式\cref{eq:1}の$2$解はこれらである．
故に解と係数の関係から
\begin{align}
  \begin{cases}
    p\equiv\cos\left(\dfrac{2\pi}{5}\right)+\cos\left(\dfrac{4\pi}{5}\right)=\dfrac{-1}{2}  \\
    q\equiv\cos\left(\dfrac{2\pi}{5}\right)\cos\left(\dfrac{4\pi}{5}\right)=\dfrac{-1}{4}
  \end{cases} \label{eq:2}
\end{align}
である．

\subparagraph{(1)}

倍角公式と\cref{eq:2}から
\begin{align}
  \begin{cases}
    a = \dfrac{1}{2}\left(1-\cos\left(\dfrac{2\pi}{5}\right)\right) \\
    b = \dfrac{1}{2}\left(1-\cos\left(\dfrac{4\pi}{5}\right)\right)
  \end{cases}
\end{align}
である．よって題意の$a+b,ab$は
\begin{align}
  \begin{cases}
    a + b & = \dfrac{1}{2}\left(1-\cos\left(\dfrac{2\pi}{5}\right)\right) + \dfrac{1}{2}\left(1-\cos\left(\dfrac{4\pi}{5}\right)\right) = \dfrac{1}{2}(2-p) = \dfrac{4}{5} \\
    ab    & = \dfrac{1}{2}\left(1-\cos\left(\dfrac{2\pi}{5}\right)\right)\cdot\dfrac{1}{2}\left(1-\cos\left(\dfrac{4\pi}{5}\right)\right) = \dfrac{1}{4}\left(1-p+q\right) = \dfrac{5}{16}
  \end{cases} \label{eq:3}
\end{align}
である．故に$a+b$も$ab$も有理数である．
よって示された．$\Box$

\subparagraph{(2)}

新しく
\begin{align}
  \begin{cases}
    a_n=(a^{-n}+b^{-n})(a+b)^n \\
    b_n=a^{-n}+b^{-n}
  \end{cases}
\end{align}
とおく．まず$b_n$についての漸化式は
\begin{align}
  b_{n+2}=\left(\dfrac{1}{a}+\dfrac{1}{b}\right)b_{n+1}-\dfrac{1}{ab}b_n
\end{align}
である．よって$a_n=b_n(a+b)^n$に注意して
\begin{align}
  a_{n+2}
  &=b_{n+2}(a+b)^{n+2} \\
  &=\left\{(\dfrac{1}{a}+\dfrac{1}{b})b_{n+1}-\dfrac{1}{ab}b_n\right\}(a+b)^n \\
  &=\dfrac{(a+b)^2}{ab}(a_{n+1}-a_n) \\
  &=5(a_{n+1}-a_n) \quad (\because\text{\cref{eq:3}})
\end{align}
である．これと$a_1=5$，$a_2=15$から，帰納的に$a_n\in\mathbb{Z}$である．$\Box$
\newpage
\end{document}
